\section{Uniform correction to the actual Ricci-flat metrics}
\label{sec:ricci-flat}

We now solve the Monge--Ampère equation on the prescribed fibres
in their prescribed polarization. The six quadratic harmonic
directions at each node are represented by finite automorphisms,
so that their nonlinear volume contribution cancels exactly.

\begin{theorem}\label{thm:ricci-flat-comparison}
Let $\omega_p$ be the Ricci-flat Kähler form in $c_1(\mathcal L_p)$.
With the fixed choices of Section~\ref{sec:global-inverse}, there
are $R_1>0$ and $C_k$ such that, for each fixed $0<R<R_1$
and sufficiently small $p\in B^\circ$,
\[
 \omega_p=\widetilde\omega_p^R
                     +\mathrm i\partial\bar\partial u_p^R,
 \qquad Q(u_p^R)=0,
\]
where $u_p^R$ is real and
\begin{equation}\label{eq:actual-potential-bound}
 \limsup_{\substack{p\to0\\p\in B^\circ}}
 \max_{j\le k}\sup_{X_p}\varrho_p^j
 |\nabla^j(\mathrm i\partial\bar\partial u_p^R)|
                  _{\widetilde g_p^R}
       \le C_kR^{\mu-\eta}.
\end{equation}
Consequently the compatible surgery identifications satisfy
\begin{equation}\label{eq:actual-model-comparison}
 \limsup_{\substack{p\to0\\p\in B^\circ}}\sup
 |\Phi_p^*g_p-\bar g_{\varepsilon(p)}^{R,A}|
                   _{\bar g_{\varepsilon(p)}^{R,A}}
 \le C R^{\mu-\eta}+CR^\mu+C e_{\mathrm{AC}}(A).
\end{equation}
There is no restriction on scale ratios or phases. The limits
are taken with $R,A$ fixed before $p\to0$, and then
$R\downarrow0$, $A\to\infty$. No parameter-derivative estimate
or arbitrary simultaneous choice $R=R(p)$ is asserted.
\end{theorem}

\subsection{The nonlinear potential coordinates}

Use the moment decomposition of
\eqref{eq:global-domain-norm}. Separate the nine
pluriharmonic coefficients at each node from the
six quadratic coefficients, and identify the
latter with $A_i\in\mathrm i\mathfrak{so}(4,\RR)$
using the fixed basis in
Lemma~\ref{lem:cap-harmonic-basis}. For the Stenzel
potential $\psi_{i,p}$ define
\[
 \Psi_{i,p}(A_i)(z)
       =\psi_{i,p}(e^{A_i}z)-\psi_{i,p}(z).
\]
After shrinking the fixed charts this is defined
for $A_i$ in a fixed small ball. Let $E_p^R$ be
the product of the $Q$-zero, zero-core-moment
remainder space and the finite coefficients
$b=(b_{i,a})_{a\le9}$, $A=(A_i)$, with norm
\[
 \|x\|_E=\|v\|_{C^{k+2,\alpha}_{2+\eta}}
                         +\sum_{i,a}|b_{i,a}|+\sum_i|A_i|.
\]
For $x=(v,b,A)$ put
\begin{equation}\label{eq:nonlinear-potential-map}
 S_p(x)=v+\sum_i\chi_i
       \left(\sum_{a=1}^9b_{i,a}H_{i,a,p}
                             +\Psi_{i,p}(A_i)\right).
\end{equation}
The derivative $J_p=DS_p(0)$ is the canonical
linear identification with $X_{p,0}^R$ and has
uniformly bounded inverse by the norm definition.
Also $Q(S_p(x))=0$.

For small $A$, the differentiated radial Stenzel
formula gives
\begin{equation}\label{eq:finite-automorphism-bounds}
 \begin{split}
 |\nabla^j\Psi_{i,p}(A)|
       &\le C_j|A|\varrho_p^{2-j},\\
 |\nabla^j(\sigma_{i,p,A}-\sigma_{i,p})|
       &\le C_j|A|\varrho_p^{-j},
 \qquad \sigma_{i,p,A}=(e^A)^*\sigma_{i,p}.
 \end{split}
\end{equation}
Positive-order derivatives in $A$ have the same
radius weights, with products of direction norms
in place of $|A|$. The estimates include scaled
Hölder norms. On the unit cap they follow from
smoothness on its core and the differentiated
asymptotic expansion on its end; the matrices
$e^A$ and their inverses range in a compact set,
so radii and scaled derivatives are comparable.
Homothety proves uniformity in every
$\varepsilon_i$.

The exact identity
\begin{equation}\label{eq:finite-automorphism-volume}
 \sigma_{i,p,A}^{\,3}=\sigma_{i,p}^{\,3}
\end{equation}
is decisive. The matrix $e^A$ preserves the
quadric and residue form because
$A^{\mathsf T}=-A$ and $\tr A=0$. Pulling back
the normalized Stenzel Calabi--Yau equation
proves \eqref{eq:finite-automorphism-volume}
for finite $A$.

We do not assert that $S_p$ maps a uniform
parameter ball into a uniform ball of the
linear norm $X_p^R$. Its nonlinear degree-two
angular profile need not have that property.
Instead we estimate the volume operator
directly on $E_p^R$. The variable remainder
$v$ is not pulled back by $e^A$, so no
differentiability loss from composition of
variable Hölder functions is involved.

\subsection{Uniform weighted nonlinear estimates}

Define
\begin{equation}\label{eq:nonlinear-log-volume}
 \mathcal M_p^R(x)=
  \log\frac{(\widetilde\omega_p^R+
            \mathrm i\partial\bar\partial S_p(x))^3}
                 {(\widetilde\omega_p^R)^3}.
\end{equation}
On a fixed sufficiently small ball in $E_p^R$
the perturbed form is positive,
$\mathcal M_p^R(x)\in Y_p^R$, and
\begin{equation}\label{eq:nonlinear-derivative-bound}
 \|D\mathcal M_p^R(x)-D\mathcal M_p^R(y)\|_{E\to Y}
       \le C\|x-y\|_E.
\end{equation}
Here and below the radius and constants are
independent of $R,p$.

We verify the weight near a node. Write
$e_i=\widetilde\omega_p^R-\sigma_{i,p}$;
\eqref{eq:initial-local-derivatives} gives
$e_i=O(\varrho_p^\mu)$ with scaled derivatives.
On $r_i<r_0$ the cutoffs are constant and the
pluriharmonic terms have zero complex Hessian.
The perturbed form is therefore
\[
 \sigma_{i,p,A_i}+e_i+B,\qquad
 B=\mathrm i\partial\bar\partial v
       =O(\|v\|_E\varrho_p^\eta).
\]
By \eqref{eq:finite-automorphism-volume}, its
logarithmic volume ratio is exactly
\begin{equation}\label{eq:weighted-logdet-cancellation}
 \log\det\bigl(I+\sigma_{i,p,A_i}^{-1}(e_i+B)\bigr)
 -\log\det\bigl(I+\sigma_{i,p}^{-1}e_i\bigr).
\end{equation}
For small fixed $r_0$ and parameter radius,
these matrices remain in a fixed positive
cone. Inversion and logarithmic determinant
have bounded derivatives there. Each variation
in $v$ has weight $\eta$. A variation in $A$
differentiates $\sigma_A^{-1}$ but retains
the factor $e_i+B$; all pure $A$ derivatives
vanish when $e_i=B=0$. Two $v$ variations have
weight $2\eta$, and mixed variations have
weight at least $\eta$, since $\mu>\eta$.
Scaled Hölder product estimates prove membership
in $Y$ and \eqref{eq:nonlinear-derivative-bound}.

On the cutoff annuli and exterior, radii are
bounded below by fixed $r_0$, so ordinary
Hölder matrix estimates apply. Globally,
\eqref{eq:finite-automorphism-bounds} gives
\begin{equation}\label{eq:nonlinear-tensor-bound}
 \max_{j\le k}\sup\varrho_p^j
 |\nabla^j(\mathrm i\partial\bar\partial S_p(x))|
                    _{\widetilde g_p^R}\le C\|x\|_E.
\end{equation}
Shrinking the parameter ball gives positivity.
These arguments give a $C^2$ map and the
stated uniform Lipschitz derivative.

\subsection{The global volume constraint}

Choose a nonnegative exterior bump supported
away from the $\chi_i$ and positive on an
open set. Normalize it to
$\int a_p\,\dd V_{\widetilde g_p^R}=1$.
Its $Y$ norm is uniformly bounded. The
projection
\[
 \Pi_p f=f-a_p\int_{X_p}f\,\dd V_{\widetilde g_p^R}
\]
is bounded from $Y_p^R$ to $Y_{p,0}^R$,
because the total volume is bounded and
the weight is positive. Let
$\mathcal N_p=\Pi_p\mathcal M_p^R$. Its
derivative at zero is
\[
 D\mathcal N_p(0)=P_p^R J_p:
                  E_p^R\longrightarrow Y_{p,0}^R.
\]
Theorem~\ref{thm:global-scalar-inverse} gives
a uniform inverse. The contraction
\begin{equation}\label{eq:nonlinear-contraction}
 x\longmapsto (D\mathcal N_p(0))^{-1}
 \left[\Pi_pF_p^R-
       \bigl(\mathcal N_p(x)-D\mathcal N_p(0)x\bigr)\right]
\end{equation}
therefore solves
$\mathcal N_p(x)=\Pi_pF_p^R$ for sufficiently
small $\|F_p^R\|_Y$, with a threshold independent
of $R,p$. More explicitly, the ball of radius
$2C_{\mathrm{inv}}\|\Pi_pF_p^R\|_Y$ is
preserved, and its Lipschitz constant is at
most one half once this radius is below the
fixed threshold from
\eqref{eq:nonlinear-derivative-bound}. Hence
\begin{equation}\label{eq:nonlinear-solution-bound}
 \|x_p^R\|_E\le C\|F_p^R\|_Y.
\end{equation}

The projected equation implies
$\mathcal M_p^R(x_p^R)-F_p^R=c a_p$.
Both unprojected volume integrals equal $V$:
one by \eqref{eq:initial-volume-compatibility},
the other because $S_p(x_p^R)$ is a global
potential in the same Kähler class.
The function
\[
 c\longmapsto
 \int_{X_p}e^{F_p^R+c a_p}(\widetilde\omega_p^R)^3
\]
is strictly increasing, since $a_p$ is
nonnegative and positive on a set of positive
volume. Thus $c=0$. We have solved the full
Monge--Ampère equation with its single global
compatibility condition.

\subsection{Smallness of the actual residual and uniqueness}

Proposition~\ref{prop:initial-kahler} gives
the stronger weighted residual estimate
\begin{equation}\label{eq:weighted-actual-residual}
 \limsup_{p\to0}\|F_p^R\|_{C^{k,\alpha}_\eta}
       \le C_kR^{\mu-\eta}.
\end{equation}
On the core, \eqref{eq:pluriharmonic-residual}
gives $R^{3/2-\eta}$. On the radius-$R$
transition its scaled derivatives are
$O(R^\mu)$ and the weight contributes
$R^{-\eta}$. On the fixed exterior the
residual tends smoothly to zero.
One further derivative gives the Hölder
bound. Since $\mu\le3/2$, these estimates
prove \eqref{eq:weighted-actual-residual}
on the full admissible domain.

Choose $R$ small enough for the uniform
contraction threshold, and then $p$ small
at that fixed $R$. Set $u_p^R=S_p(x_p^R)$.
Equations \eqref{eq:nonlinear-tensor-bound},
\eqref{eq:nonlinear-solution-bound} and
\eqref{eq:weighted-actual-residual} prove
\eqref{eq:actual-potential-bound}. The resulting
positive form has volume $\kappa(p)\Theta_p$,
so is Ricci-flat. Its class is exactly
$c_1(\mathcal L_p)$. Uniqueness in the fixed
Kähler class identifies it with $\omega_p$.
Elliptic bootstrapping makes it smooth, and
repeating the argument at any higher fixed
order estimates the same metric by uniqueness.

Combining \eqref{eq:actual-potential-bound}
with \eqref{eq:initial-model-comparison}
proves \eqref{eq:actual-model-comparison}.
It covers every cap, intervening annulus
and exterior point. The actual metric is
independent of the auxiliary radius $R$.
Taking successively smaller fixed $R$
also proves smooth convergence to $g_0$
on every fixed exterior compact set and
to the normalized unit Stenzel metric on
every fixed rescaled cap compact set.
This completes the proof of
Theorem~\ref{thm:ricci-flat-comparison}.
